{\small\opmaakproza
\begin{enumerate}[leftmargin=1.5em,itemsep=3pt,label=\textcolor{kleur1}{\sffamily\bfseries\arabic*}]
\item \textbf{Lees de inleiding van een opgave goed.} Gegevens die je bij een latere vraag nodig hebt, staan vaak
  al bovenaan de opgave, in een tekst, tabel of figuur.
\item \textbf{Let op het werkwoord.} Bij \emph{noem} geef je alleen het antwoord; bij \emph{leg uit} of \emph{verklaar}
  geef je ook de reden; bij \emph{bereken} schrijf je de berekening op.
\item \textbf{Geef niet meer antwoorden dan gevraagd.} Vraagt een vraag om twee redenen, dan tellen alleen de eerste twee.
\item \textbf{Gebruik de woorden uit de vraag.} Een antwoord dat duidelijk aansluit bij de vraag en de context
  scoort beter dan een algemeen verhaal.
\item \textbf{Schrijf berekeningen volledig op}, met eenheden. Bij de meeste vragen krijg je ook punten voor de tussenstappen.
\item \textbf{Kijk naar de punten.} Een vraag van 3 punten vraagt meestal drie stappen of argumenten.
  Houd ongeveer 2 minuten per punt aan en ga door als je vastloopt.
\item \textbf{Gebruik je informatieboek} als dat bij dit vak is toegestaan, en noem waar je een waarde vandaan haalt.
\item \textbf{Teken of vul in op de uitwerkbijlage} als de vraag dat zegt, en lever die in.
\item \textbf{Kom terug op vragen die je overslaat.} Zet een streepje in de kantlijn, zodat je ze aan het eind snel terugvindt.
\item \textbf{Controleer je antwoorden}: past het antwoord bij de vraag, en is het logisch in de context?
\end{enumerate}\par}
